\section{Freedom or Compulsion}

(E. E. Gynild.)

——

We need not ask which is Christian, freedom or compulsion. We have therefore sought to live in freedom in our church life, so that both contributions and services should be brought forth by inward compulsion. But freedom is dearly bought and is not preserved without struggle, struggle against the flesh, which must always be subdued, because it is an enemy of freedom. “Where the Spirit of the Lord is, there is freedom.” If the Spirit is allowed to rule, we come into freedom, as we do in the word of the Lord. Then we will not merely demand, but first and foremost give.

When I have thought about the collection for our school for pastors, this question of freedom has pressed itself upon me. Augsburg is a free institution, placed in the service of the free congregation. It has lived and preserved its freedom through many struggles and much poverty, because there were free men and women, and some free congregations, who willingly sacrificed and labored for its continued existence.

Freedom does not ask what others are doing. It asks only what it itself has strength to do. And in labor of that kind there is both inspiration and encouragement. It is freedom of this kind that has carried Augsburg through the times of distress through which it has passed.

In the seventies the school had a debt of \$16,000.00. The grasshoppers had ravaged the West, and our people in the new settlements were poor; nevertheless they were able to raise the money for the debt and save Augsburg from ruin. In 1917 we had a debt of \$17,000.00, and now our people are for the most part prosperous. If we now do not have the strength to make the school debt-free, then it is a very strong testimony to our lack of longing for freedom and of spirit.

I shall also remind you that in several collections since the Lutheran Free Church was organized around Augsburg, poor pastors gave subscriptions ranging from \$25 up to \$100, and among a portion of the congregational people the proportion was the same. We willed it, and we carried out what we willed. Our financial strength is greater now, but we must ask: do we will anything?

It is impossible to conceive that we should not make Augsburg debt-free this year. It would be far too sad a testimony to spiritual poverty.

We shall not complain about the collection this year where it has been carried out. But not one third of the congregations had completed their collection or sent in their contributions by New Year. Up to January 4, only 100 congregations had finished their collection. When we take into consideration the amount for which receipts had been issued, that was indeed rather encouraging. If those who remain would now do something similar, there would indeed be reason for hope. But the best time of the year for a collection has, after all, passed, and congregations in that region of the West where the crop failed can hardly do what they would have done in a good year; thus the outlook may give occasion both for fear and for hope.

This is written in order that we may settle it as our firm determination that the debt must be raised, and raised this year, so that the next fiscal year may begin without debt. I shall not weary you with much writing hereafter. Not words, but action alone will do it.

One of the older pastors writes in a letter of January 16:

\begin{quote}May the friends now in their great majority see it in this way, as he sees it, and there might come something of a mighty exertion in the work, and without doubt it would succeed.

“We are trying to struggle our way forward to faith in a successful outcome.”

“Believe, and you shall see the glory of God.”

\end{quote}

Thus far the letter.

May there arise a struggle in all believing people throughout the congregations, a struggle against worldliness and all spiritual superficiality, so that congregational life, created by the Spirit of God, may not be choked by worldliness and dead churchliness.